\documentclass[12pt]{amsart}
\usepackage{amsmath, amsthm, amssymb}
\usepackage[margin=1.2in]{geometry}
\usepackage{hyperref}
\setlength{\textheight}{8.5in}
\setlength{\evensidemargin}{0.0in}
\setlength{\oddsidemargin}{0.0in}
\setlength{\topmargin}{-0.5in}
\setlength{\textwidth}{6.5in}

\newtheorem{theorem}{Theorem}[section]
\newtheorem{lemma}[theorem]{Lemma}
\newtheorem{corollary}[theorem]{Corollary}
\newtheorem{proposition}[theorem]{Proposition}
\newtheorem{conjecture}[theorem]{Conjecture}

\theoremstyle{definition}
\newtheorem{definition}[theorem]{Definition}
\newtheorem{example}[theorem]{Example}

\theoremstyle{remark}
\newtheorem{remark}[theorem]{Remark}

\newcommand{\Z}{\mathbb{Z}}
\newcommand{\N}{\mathbb{N}}
\newcommand{\ord}{\operatorname{ord}}
\newcommand{\lcm}{\operatorname{lcm}}
\newcommand{\ind}{\operatorname{ind}}

\title{Decimal Expansions of Rational Numbers}

\begin{document}
\maketitle
\section*{Introduction}

We define the period of a fraction precisely as follows.

\begin{definition}\label{def:period}
Let $a/b$ be a rational number. The period of $a/b$ is the
smallest positive integer $k$ such that the decimal expansion of $a/b$
is eventually periodic with block length $k$. The pre-period is
the length of the non-repeating initial segment of the expansion.
\end{definition}

Definition~\ref{def:period} is not yet known to be well-posed: it
asserts the existence of a smallest $k$ for which the decimal is
periodic with that block length, but we have not established that any
such $k$ exists at all. Before proceeding, we close this gap with two
short propositions whose proofs rest on the pigeonhole principle
applied to the long-division algorithm. Both arguments are based on
the following object, which records the state of long division after
each step.

\begin{definition}\label{def:remainder-sequence}
For integers $a$ and $b$ with $b \ge 2$, the
\emph{long-division remainder sequence} of $a/b$ is the sequence
$(r_j)_{j \ge 0}$ defined by
\[
  r_0 = a \bmod b, \qquad r_{j+1} = 10\, r_j \bmod b
  \quad \text{for all } j \ge 0.
\]
\end{definition}

The remainder sequence is exactly what grade-school long division
produces: $r_j$ is the remainder carried into the $(j+1)$-st step, and
the next digit of the decimal expansion is $d_{j+1} =
\lfloor 10\, r_j / b \rfloor$. In particular, two consecutive digits
$d_{j+1}$ and $d_{j'+1}$ are equal whenever $r_j = r_{j'}$.

\begin{proposition}\label{prop:periodic-existence}
For every rational $a/b$ with $b \ge 1$, the decimal expansion of $a/b$
is eventually periodic. In particular, the period of $a/b$ in the
sense of Definition~\ref{def:period} exists.
\end{proposition}

\begin{proof}
If $b = 1$ the expansion terminates, so assume $b \ge 2$ and consider
the long-division remainder sequence $(r_j)_{j \ge 0}$ from
Definition~\ref{def:remainder-sequence}. Each $r_j$ takes a value in
the finite set $\{0, 1, \ldots, b - 1\}$, which has $b$ elements. Among
the first $b + 1$ terms $r_0, r_1, \ldots, r_b$, the pigeonhole
principle guarantees that two must coincide: $r_i = r_j$ for some
$0 \le i < j \le b$. The recursion $r_{k+1} = 10 r_k \bmod b$ depends
only on $r_k$, so $r_i = r_j$ forces $r_{i+t} = r_{j+t}$ for every
$t \ge 0$. Since each digit $d_{k+1}$ is determined by $r_k$, the
digits from position $i + 1$ onward repeat with block length at most
$j - i \le b$. Thus the decimal expansion is eventually periodic, and
Definition~\ref{def:period} yields a well-defined smallest $k$.
\end{proof}

A more careful application of pigeonhole — restricted to remainders
that are coprime to $b$ — yields a quantitative bound involving
Euler's totient function $\varphi$, defined by
$\varphi(b) = |(\Z/b\Z)^*|$.

\begin{proposition}\label{prop:phi-bound}
Let $a/b$ be in lowest terms with $\gcd(b, 10) = 1$ and $b \ge 2$.
Then the period of $a/b$ is at most $\varphi(b)$.
\end{proposition}

\begin{proof}
We first show that every term of the long-division remainder sequence
$(r_j)$ lies in the multiplicative group $(\Z/b\Z)^*$, i.e.\ is
coprime to $b$. Since $a/b$ is in lowest terms, $\gcd(a, b) = 1$, and
reducing $a$ modulo $b$ preserves coprimality, so $\gcd(r_0, b) = 1$.
Suppose inductively that $\gcd(r_j, b) = 1$. Because
$\gcd(10, b) = 1$, we have $\gcd(10\, r_j, b) = 1$, and reducing
$10\, r_j$ modulo $b$ once again preserves coprimality, giving
$\gcd(r_{j+1}, b) = 1$. By induction, every $r_j$ is a unit modulo
$b$.

The sequence $(r_j)$ therefore takes values in the set $(\Z/b\Z)^*$,
which has exactly $\varphi(b)$ elements. Among the first
$\varphi(b) + 1$ terms, the pigeonhole principle forces two to
coincide, say $r_i = r_j$ with $0 \le i < j \le \varphi(b)$. The same
argument as in the proof of Proposition~\ref{prop:periodic-existence}
shows that the digit sequence from position $i + 1$ onward repeats
with block length at most $j - i \le \varphi(b)$, so the period is at
most $\varphi(b)$.
\end{proof}

\begin{example}\label{ex:phi-not-tight}
Proposition~\ref{prop:phi-bound} is rarely tight. For $b = 13$,
$\varphi(13) = 12$, but
\[
  \frac{1}{13} = 0.\overline{076923}
\]
has period $6$, exactly half the bound. For $b = 21 = 3 \cdot 7$,
$\varphi(21) = 12$, yet
\[
  \frac{1}{21} = 0.\overline{047619}
\]
also has period $6$. In both cases the true period is a proper divisor
of $\varphi(b)$.
\end{example}

Proposition~\ref{prop:phi-bound} gives the right order of magnitude
but can be loose by a factor of $2$ or more, and offers no insight
into \emph{which} divisor of $\varphi(b)$ the period equals. The goal
of this paper is to give a complete formula that
takes the prime factorization of $b$ as input and returns the period and
pre-period of any fraction in the lowest terms $a/b$. The result is the following.

\medskip
\begin{theorem}\label{lem:main-formula}
Let $a/b$ be a fraction in lowest terms with
\[
  b = 2^\alpha \cdot 5^\beta \cdot p_1^{e_1} \cdots p_r^{e_r},
\]
where the $p_i$ are distinct primes coprime to $10$. Then the decimal
expansion of $a/b$ has exactly $L = \max(\alpha, \beta)$ non-repeating
digits, followed by a purely repeating block of length
\[
  \lcm\!\bigl(\lambda(p_1^{e_1}),\, \ldots,\, \lambda(p_r^{e_r})\bigr),
\]
where each
\[
  \lambda(p_i^{e_i}) \;=\; \lambda(p_i)\cdot p_i^{\max(e_i - 1 - s_i,\; 0)},
  \qquad s_i = \nu_{p_i}(10^{\lambda(p_i)} - 1) - 1.
\]
\end{theorem}
\medskip

\section{The Period as a Multiplicative Order}\label{sec:order}

Our first task is to reduce the definition of the period, a statement about decimal digits, into a statement about divisibility, and then into a statement about multiplicative order in $\Z / b \Z$. We carry out this reduction 
in two steps. The first step, valid for any numerator, reformulates the 
period as the smallest $k$ for which $b \mid a(10^k - 1)$. The second 
step uses the coprimality of $a$ and $b$ to further to a
condition on $10^k$ alone. This reveals that the period of $a/b$ is 
equal to $\ord_b(10)$, and in particular does not depend on $a$.

Before connecting the period to a divisibility condition, we note that 
the assumption $\gcd(b, 10) = 1$ rules out any non-repeating prefix: 
the decimal expansion of $a/b$ is purely periodic, with the repeating 
block beginning immediately after the decimal point.

\begin{lemma}\label{lem:purely-periodic}
Let $a/b$ be a fraction in lowest terms with $\gcd(b, 10) = 1$ and 
$b \ge 2$. Then the decimal expansion of $a/b$ is purely periodic with pre-period length 0.
\end{lemma}

\begin{proof}
Suppose for contradiction that the expansion of $a/b$ has a 
non-repeating prefix of length $s \ge 1$, followed by a repeating block 
of some length $k$. So we may write
\[
  \frac{a}{b} = 0.d_1 d_2 \cdots d_s \,\overline{d_{s+1} \cdots d_{s+k}},
\]
where the digits $d_1, \ldots, d_s$ form the non-repeating prefix and 
$d_{s+1} \cdots d_{s+k}$ is the block that repeats forever.

Consider what happens when we multiply the fraction by $10^s$. In general, 
multiplying a decimal by $10^s$ shifts the decimal point $s$ places to 
the right, so the digits to the left of the new decimal point become the 
integer part, and everything to the right becomes the new fractional 
part. Applied to $a/b$,
\[
  10^s \cdot \frac{a}{b} 
  = d_1 d_2 \cdots d_s . \,\overline{d_{s+1} \cdots d_{s+k}}.
\]
The integer part is $d_1 d_2 \cdots d_s$, and the fractional part 
consumes exactly the non-repeating prefix: what remains after the new 
decimal point begins with $d_{s+1}$, which by assumption is the start 
of the repeating block. So the fractional part of $10^s \cdot (a/b)$ is 
purely periodic with period $k$. 

Now, a number with a purely periodic decimal of period $k$ is fixed in 
its fractional part when we shift by a further $k$ places: multiplying 
by $10^k$ and subtracting the original leaves an integer, because the 
fractional parts cancel. Applying this to $10^s \cdot (a/b)$:
\[
  (10^k - 1) \cdot 10^s \cdot \frac{a}{b} 
  = 10^{s+k} \cdot \frac{a}{b} - 10^s \cdot \frac{a}{b} \in \Z.
\]
From this, we see that $b \mid 10^s(10^k - 1)\,a$. Since 
$\gcd(a, b) = 1$, we may cancel $a$ to obtain $b \mid 10^s(10^k - 1)$, 
and since $\gcd(b, 10) = 1$ gives $\gcd(b, 10^s) = 1$, we conclude 
$b \mid (10^k - 1)$.

However, from this last statement we can derive $(10^k - 1)(a/b) \in \Z$, thus, multiplying $a/b$ itself (no preliminary shift by $10^s$) by 
$10^k - 1$ gives an integer. Using our previous argument's logic in 
reverse, this means the fractional part of $a/b$ is already purely 
periodic with period $k$, with no non-repeating prefix at all. This 
contradicts the assumption $s \ge 1$.
\end{proof}

Now that we have a condition for pure periodicity in hand, we can now translate the period into a 
divisibility condition.

\subsection{Connecting the period to a congruence condition}

The following lemma is the bridge between decimal expansions and 
congruence conditions. It follows from the observation that shifting the decimal 
point by $k$ positions corresponds to multiplication by $10^k$, and that for a fraction that has period of length $k$, this shift
preserves the repeating fractional part.

\begin{lemma}\label{lem:period-as-divisibility}
Let $a/b$ be a fraction with $\gcd(b, 10) = 1$. The period of $a/b$ is 
the smallest positive integer $k$ such that $b \mid a(10^k - 1)$.
\end{lemma}

\begin{proof}
The decimal expansion of $a/b$ repeats with block length $k$ if and only 
if shifting by $k$ decimal places reproduces the fractional part as such:
\[
  10^k \cdot \frac{a}{b} - \frac{a}{b} \in \Z.
\]
Factoring the left-hand side yields $(10^k - 1)\,a/b \in \Z$, which 
holds if and only if $b \mid a(10^k - 1)$. The period is by definition 
the smallest $k$ for which this holds.
\end{proof}

\subsection{The multiplicative order and canceling the numerator}

To remove the numerator from the condition $b \mid a(10^k - 1)$ we 
need a basic fact: when $a$ and $b$ are coprime, $a$ is a unit in 
$\Z/b\Z$ and can be cancelled from any divisibility statement in which 
it appears as a factor. Before formalizing this, we introduce the following definition that will be used throughout the proofs of this paper.

\begin{definition}\label{def:order}
For integers $n$ and $m$ with $\gcd(n, m) = 1$ and $m \ge 2$, the 
\emph{multiplicative order} of $n$ modulo $m$, written $\ord_m(n)$, 
is the smallest positive integer $k$ such that $n^k \equiv 1 \pmod{m}$.
\end{definition}

That $\ord_m(10)$ exists whenever $\gcd(10, m) = 1$ requires brief 
justification. The condition $\gcd(10, m) = 1$ says that $10$ is a unit 
in $\Z/m\Z$, and hence belongs to the multiplicative group 
$(\Z/m\Z)^*$, which is finite of order $\varphi(m)$. Every element of a 
finite group has finite order, so there is some positive integer $k$ 
with $10^k \equiv 1 \pmod{m}$, and the smallest such $k$ is well-defined. 
We note that this argument applies to any $m$ coprime to $10$, not only 
prime values.

The following divisibility property follows from Definition~\ref{def:order} and 
will be invoked repeatedly in Sections~\ref{sec:reduction} 
and~\ref{sec:prime-power}.

\begin{lemma}\label{lem:order-divides}
If $10^k \equiv 1 \pmod{m}$, then $\ord_m(10) \mid k$.
\end{lemma}

\begin{proof}
Let $d = \ord_m(10)$. By the division algorithm, write $k = qd + r$ with 
$0 \le r < d$. Then
\[
  10^r = 10^{k - qd} = 10^k \cdot (10^d)^{-q} \equiv 1 \cdot 1^{-q} = 1 
  \pmod{m}.
\]
If $r > 0$, this would contradict the minimality of $d$, since $r$ 
would be a smaller positive integer satisfying the same congruence. 
Hence $r = 0$, thus, $d \mid k$.
\end{proof}

We are now ready to combine the two steps. 
Lemma~\ref{lem:period-as-divisibility} reformulates the period as a 
divisibility condition, and the coprimality of $a$ and $b$ allows us to reduce to a pure congruence on $10^k$. The minimal such $k$ is 
then $\ord_b(10)$.

\begin{theorem}\label{thm:period-equals-order}
Let $a/b$ be in lowest terms with $\gcd(b, 10) = 1$ and $b \ge 2$. 
Then the period of $a/b$ equals $\ord_b(10)$. In particular, the period 
depends only on $b$, not on the numerator $a$.
\end{theorem}

\begin{proof}
By Lemma~\ref{lem:period-as-divisibility}, the period is the smallest 
positive integer $k$ such that $b \mid a(10^k - 1)$. Since 
$\gcd(a, b) = 1$, the residue of $a$ is a unit in $\Z/b\Z$, so
\[
  b \mid a(10^k - 1) 
  \;\;\Longleftrightarrow\;\; b \mid (10^k - 1) 
  \;\;\Longleftrightarrow\;\; 10^k \equiv 1 \pmod{b}.
\]
The smallest positive $k$ satisfying the right-hand condition is 
$\ord_b(10)$, which depends only on $b$.
\end{proof}

\begin{corollary}\label{cor:numerator-irrelevant}
If $a, a'$ are both coprime to $b$ and $\gcd(b, 10) = 1$, then $a/b$ 
and $a'/b$ have the same period.
\end{corollary}

This corollary explains the coincidence noted in the introduction: $2/7$, 
$3/7$, and $6/7$ all have period $6$ because the period depends only on 
the denominator. Theorem~\ref{thm:period-equals-order} requires 
$\gcd(b, 10) = 1$; the next section handles the general case.

\section{Reducing the Period of Any Fraction to Prime Powers}
\label{sec:reduction}

We now remove the restriction $\gcd(b, 10) = 1$ and allow arbitrary 
denominators. The result is a two-stage reduction: factors of $2$ and 
$5$ in $b$ produce a non-repeating prefix, and the remaining coprime-to-$10$ 
part $m$ determines the periodic portion. Moreover, if $m$ is composite, 
the Chinese Remainder Theorem further reduces $\ord_m(10)$ to a least 
common multiple of prime-power orders, so that the full computation is 
controlled by the values $\ord_{p^e}(10)$ for prime powers $p^e$ with 
$p \nmid 10$. Those values are the subject of Section~\ref{sec:prime-power}.

\subsection{The pre-period: stripping out factors of 2 and 5}

Throughout the remainder of the paper we adopt the following notation. 
For any integer $b \ge 2$, write
\[
  b = 2^\alpha \cdot 5^\beta \cdot m, \qquad \gcd(m, 10) = 1,
\]
where $\alpha = \nu_2(b)$, $\beta = \nu_5(b)$, and $m$ is the 
coprime-to-$10$ part of $b$. Define $L = \max(\alpha, \beta)$.

\begin{theorem}\label{thm:preperiod}
Let $a/b$ be in lowest terms with $b = 2^\alpha \cdot 5^\beta \cdot m$ 
and $\gcd(m, 10) = 1$. Then the decimal expansion of $a/b$ has exactly 
$L = \max(\alpha, \beta)$ non-repeating digits. If $m = 1$ the expansion 
terminates after $L$ digits; if $m > 1$, these $L$ digits are followed 
by a purely periodic block of length $\ord_m(10)$.
\end{theorem}

\begin{proof}
\textbf{Upper bound ($L$ steps suffice).} Compute
\[
  10^L \cdot \frac{a}{b} 
  = \frac{2^L \cdot 5^L \cdot a}{2^\alpha \cdot 5^\beta \cdot m}
  = \frac{a \cdot 2^{L - \alpha} \cdot 5^{L - \beta}}{m}.
\]
Both exponents $L - \alpha$ and $L - \beta$ are non-negative by the 
definition of $L$, so the numerator $a' = a \cdot 2^{L-\alpha} \cdot 
5^{L-\beta}$ is a positive integer. The denominator is $m$, which is 
coprime to $10$. Moreover $\gcd(a', m) = 1$: since $a/b$ is in lowest 
terms we have $\gcd(a, m) = 1$ (because $m \mid b$), and 
$\gcd(2, m) = \gcd(5, m) = 1$ by the definition of $m$, so the extra 
factors of $2$ and $5$ introduce no common factor with $m$. Thus 
$a'/m$ is in lowest terms with denominator coprime to $10$. If $m = 1$, 
the fraction $a'/m$ is an integer and the decimal of $a/b$ terminates; 
if $m > 1$, Theorem~\ref{thm:period-equals-order} gives that $a'/m$ is 
purely periodic with period $\ord_m(10)$. In either case, multiplying 
by $10^L$ shifts the decimal point $L$ places to the right and produces 
a number whose fractional part is purely periodic, so the pre-period of 
$a/b$ is at most $L$.

\smallskip

\textbf{Lower bound (fewer than $L$ steps cannot work).} Suppose for 
contradiction that the expansion of $a/b$ becomes purely periodic 
starting at position $s + 1$ for some $s < L$. Then $10^s \cdot (a/b)$ 
has a purely periodic fractional part with some period $k$, which means
\[
  (10^k - 1) \cdot 10^s \cdot \frac{a}{b} \in \Z,
\]
so $b \mid 10^s (10^k - 1)\, a$. Since $\gcd(a, b) = 1$, we may cancel 
$a$ to obtain $b \mid 10^s(10^k - 1)$. In particular $2^\alpha \mid 
10^s(10^k - 1)$; but $10^k - 1$ is odd (every power of $10$ ends in $0$, 
so $10^k - 1$ ends in $9$), so $\nu_2(10^k - 1) = 0$, and all 
$\alpha$ factors of $2$ must come from $10^s$. This forces $s \ge 
\alpha$. An analogous argument can be made with the prime $5$: since $10 \equiv 0 
\pmod{5}$, we have $10^k - 1 \equiv -1 \equiv 4 \pmod{5}$, so 
$5 \nmid (10^k - 1)$, and thus $s \ge \beta$. Combining, 
$s \ge \max(\alpha, \beta) = L$, contradicting $s < L$. 

Therefore the pre-period is at least $L$, and combined with the upper 
bound just proved, it must be exactly $L$.
\end{proof}

We illustrate the theorem with the following example.

\begin{example}\label{ex:1-over-56}
For $1/56$, we have $b = 56 = 2^3 \cdot 7$, so $\alpha = 3$, $\beta = 0$, 
$m = 7$, and $L = 3$. The theorem predicts three non-repeating digits 
followed by a block of length $\ord_7(10) = 6$. Indeed, 
$1/56 = 0.017\overline{857142}$.
\end{example}

\subsection{Decomposing $\ord_m(10)$ over prime powers}

Theorem~\ref{thm:preperiod} reduces the computation of periods to 
understanding $\ord_m(10)$ for integers $m$ coprime to $10$. When 
$m$ is itself a prime power this is a single quantity, computed 
explicitly in Section~\ref{sec:prime-power}. When $m$ is composite, 
the Chinese Remainder Theorem reduces the problem to its prime-power 
factors.

\begin{lemma}\label{lem:ord-lcm}
Let $m = p_1^{e_1} \cdots p_r^{e_r}$ be a factorization into distinct 
prime powers, each coprime to $10$. Then
\[
  \ord_m(10) = \lcm\!\bigl(\ord_{p_1^{e_1}}(10),\, \ldots,\, 
  \ord_{p_r^{e_r}}(10)\bigr).
\]
\end{lemma}

\begin{proof}
Let $k = \ord_m(10)$ and $L = \lcm(\ord_{p_1^{e_1}}(10), \ldots, 
\ord_{p_r^{e_r}}(10))$.

To show $k \mid L$: each $\ord_{p_i^{e_i}}(10)$ divides $L$ by 
definition of $\lcm$, so $10^L \equiv 1 \pmod{p_i^{e_i}}$ for every 
$i$. We want to combine these $r$ separate congruences into a single 
congruence modulo $m = p_1^{e_1} \cdots p_r^{e_r}$. This is exactly 
what the Chinese Remainder Theorem lets us do: when the moduli 
$p_1^{e_1}, \ldots, p_r^{e_r}$ are pairwise coprime, an integer $N$ 
satisfies $N \equiv 0 \pmod{p_i^{e_i}}$ for every $i$ if and only if 
$N \equiv 0 \pmod{m}$. Applying this with $N = 10^L - 1$, the 
congruences $10^L \equiv 1 \pmod{p_i^{e_i}}$ for all $i$ combine to 
give $10^L \equiv 1 \pmod{m}$. By Lemma~\ref{lem:order-divides}, 
$k \mid L$.

To show $L \mid k$: $10^k \equiv 1 \pmod{m}$ implies $10^k \equiv 1 
\pmod{p_i^{e_i}}$ for every $i$, so by Lemma~\ref{lem:order-divides} 
each $\ord_{p_i^{e_i}}(10)$ divides $k$. Their least common multiple 
$L$ then also divides $k$.

Since $k \mid L$ and $L \mid k$ are both positive integers, $k = L$.
\end{proof}

Lemma~\ref{lem:ord-lcm} reduces the computation of the period for any 
fraction to the computation of $\ord_{p^e}(10)$ for each prime power 
$p^e$ in the factorization of the denominator, which will be discussed in Section~\ref{sec:prime-power}.

\begin{example}\label{ex:1-over-42}
For $1/42$, we have $b = 42 = 2 \cdot 3 \cdot 7$, so $\alpha = 1$, 
$\beta = 0$, $m = 21$, and $L = 1$. By Lemma~\ref{lem:ord-lcm},
\[
  \ord_{21}(10) = \lcm(\ord_3(10), \ord_7(10)) = \lcm(1, 6) = 6,
\]
taking the prime-power values as given here; 
Section~\ref{sec:prime-power} justifies them. So $1/42$ has one 
non-repeating digit followed by a block of length $6$. Indeed, 
$1/42 = 0.0\overline{238095}$.
\end{example}


%%-------------------------------------------------
\section{The Period of $1/p^n$}\label{sec:prime-power}
%%-------------------------------------------------

The remaining computation is that of $\ord_{p^n}(10)$ for a single 
prime power $p^n$ with $p \nmid 10$. For brevity, write 
$\lambda(m) = \ord_m(10)$ throughout this section, and set 
$\lambda = \lambda(p) = \ord_p(10)$. We also need the $p$-adic 
valuation, which is defined as follows:

\begin{definition}\label{def:p-adic}
For a prime $p$ and a nonzero integer $m$, the \emph{$p$-adic valuation} 
$\nu_p(m)$ is the largest non-negative integer $k$ such that $p^k \mid m$.
\end{definition}
Define $s = \nu_p(10^\lambda - 1) - 1$, which is non-negative since 
$p \mid (10^\lambda - 1)$ by definition of $\lambda$.
We now show the following.

\begin{theorem}\label{thm:prime-power-period}
For any prime $p$ with $p \nmid 10$, and with $\lambda = \ord_p(10)$ 
and $s = \nu_p(10^\lambda - 1) - 1$,
\[
  \lambda(p^n) = \lambda \cdot p^{\max(n - 1 - s,\, 0)} 
  \qquad \text{for all } n \ge 1.
\]
\end{theorem}

The proof proceeds in two stages. Section~\ref{subsec:form} shows that 
$\lambda(p^n)$ must take the form $\lambda \cdot p^e$ for some 
$0 \le e \le n - 1$, using Fermat's Little Theorem and a binomial 
expansion to establish divisibility in both directions. 
Section~\ref{subsec:lifting} introduces a lemma on how $\nu_p(a^p - 1)$ 
grows with iteration, and Section~\ref{subsec:pinning} combines these to 
pin down the value of $e$ exactly. 


\subsection{Proving $\lambda(p^n)$ has form $\lambda \cdot p^e$}\label{subsec:form}

We now constrain $\lambda(p^n)$ for $n \ge 1$. The argument has two 
halves: $\lambda$ divides $\lambda(p^n)$ because any period modulo 
$p^n$ is also a period modulo $p$; and $\lambda(p^n)$ divides 
$\lambda \cdot p^{n-1}$ by a standard binomial expansion.

\begin{lemma}\label{lem:form-lambda-pe}
For every $n \ge 1$, $\lambda(p^n) = \lambda \cdot p^e$ for some integer 
$0 \le e \le n - 1$.
\end{lemma}

\begin{proof}
Let $d = \lambda(p^n)$. 

\smallskip

First, we show $\lambda \mid d$. Since $10^d \equiv 1 \pmod{p^n}$ and 
$p \mid p^n$, reducing modulo $p$ gives $10^d \equiv 1 \pmod{p}$. By 
Lemma~\ref{lem:order-divides} applied at $m = p$, $\lambda \mid d$. 

We will now show that $d \mid \lambda \cdot p^{n-1}$. The strategy is 
to prove that $10^{\lambda \cdot p^{n-1}} \equiv 1 \pmod{p^n}$; 
Lemma~\ref{lem:order-divides} then gives the divisibility.

Since $p \mid (10^\lambda - 1)$ by definition of $\lambda$, we may 
write $10^\lambda = 1 + cp$ for some integer $c$. Raising both sides 
to the power $p^{n-1}$,
\[
  10^{\lambda \cdot p^{n-1}} = (1 + cp)^{p^{n-1}},
\]
and expanding the right-hand side by the binomial theorem,
\[
  (1 + cp)^{p^{n-1}} = 1 + \sum_{k=1}^{p^{n-1}} \binom{p^{n-1}}{k} 
  (cp)^k.
\]
To establish that this equals $1$ modulo $p^n$, it suffices to show 
that every term in the sum is divisible by $p^n$, which is equivalent to showing that 
$\nu_p\!\bigl(\binom{p^{n-1}}{k}(cp)^k\bigr) \ge n$ for each 
$1 \le k \le p^{n-1}$.

Each term has two sources of $p$-divisibility: the binomial coefficient 
and the power $(cp)^k$. We analyze each in turn.

For the binomial coefficient, a standard identity (derived by direct manipulation of factorials) states that 
for any prime $p$ and any $1 \le k \le p^{n-1}$,
\[
  \nu_p\!\binom{p^{n-1}}{k} = (n-1) - \nu_p(k).
\]
Intuitively: $\binom{p^{n-1}}{k}$ carries a factor of $p^{n-1}$ from 
the numerator $p^{n-1}!$, partially cancelled by the $p$-divisibility 
of $k!$ in the denominator. The $p$-adic valuation of $k!$ accounts 
for exactly $\nu_p(k)$ of those $p$'s (this is where $k$'s own 
divisibility by $p$ enters), so $n - 1 - \nu_p(k)$ factors of $p$ 
survive in the binomial coefficient.

For the power $(cp)^k$, we simply have $\nu_p\!\bigl((cp)^k\bigr) \ge 
k$, with equality when $p \nmid c$.

Combining, the $p$-adic valuation of the $k$-th term is at least
\[
  \bigl(n - 1 - \nu_p(k)\bigr) + k.
\]
We now show this is at least $n$, mainly, that $k - \nu_p(k) \ge 1$ for 
every positive integer $k$. Writing $k = p^{\nu_p(k)} \cdot k'$ with 
$\gcd(k', p) = 1$, we have $k \ge p^{\nu_p(k)} \ge \nu_p(k) + 1$, where 
the second inequality is the familiar fact that $p^j \ge j + 1$ for 
every $j \ge 0$ and every prime $p \ge 2$ (by induction: $p^{j+1} \ge 
2 p^j \ge 2(j + 1) \ge j + 2$). Therefore $k - \nu_p(k) \ge 1$, and 
the $k$-th term has $\nu_p \ge n$ as claimed.

Every term in the sum is thus divisible by $p^n$, so
\[
  10^{\lambda \cdot p^{n-1}} = 1 + \sum_{k=1}^{p^{n-1}} 
  \binom{p^{n-1}}{k}(cp)^k \equiv 1 \pmod{p^n}.
\]
By Lemma~\ref{lem:order-divides} applied at $m = p^n$, 
$d \mid \lambda \cdot p^{n-1}$.

% We will now show that $d \mid \lambda \cdot p^{n-1}$. Write $10^\lambda = 1 + cp$ 
% for some integer $c$, which is possible since $p \mid (10^\lambda - 1)$. 
% Then
% \[
%   10^{\lambda \cdot p^{n-1}} = (1 + cp)^{p^{n-1}},
% \]
% and expanding by the binomial theorem gives
% \[
%   (1 + cp)^{p^{n-1}} = 1 + \sum_{k=1}^{p^{n-1}} \binom{p^{n-1}}{k} 
%   (cp)^k.
% \]
% For $1 \le k \le p^{n-1}$, the binomial coefficient $\binom{p^{n-1}}{k}$ 
% is divisible by $p^{n-1 - \nu_p(k)}$ (a standard fact: 
% $\nu_p\!\binom{p^{n-1}}{k} = n - 1 - \nu_p(k)$ whenever $1 \le k \le 
% p^{n-1}$), and $(cp)^k$ contributes another factor of $p^k$. Thus each 
% term has $\nu_p \ge (n - 1 - \nu_p(k)) + k$, and since 
% $k \ge \nu_p(k) + 1$ for any positive integer $k$, this is at least 
% $(n - 1) + 1 = n$. Therefore every term in the sum is divisible by 
% $p^n$, and
% \[
%   10^{\lambda \cdot p^{n-1}} \equiv 1 \pmod{p^n}.
% \]
% By Lemma~\ref{lem:order-divides} applied at $m = p^n$, 
% $d \mid \lambda \cdot p^{n-1}$.

\smallskip

Combining the two conclusions, $d / \lambda$ is a positive integer 
dividing $p^{n-1}$, so $d / \lambda = p^e$ for some $0 \le e \le n - 1$.
\end{proof}

Lemma~\ref{lem:form-lambda-pe} restricts $\lambda(p^n)$ to one of the 
$n$ candidate values $\lambda, \lambda p, \lambda p^2, \ldots, 
\lambda p^{n-1}$. To determine which, we explore
how divisibility by powers of $p$ changes when we raise $10^\lambda$ 
to higher powers of $p$.

\subsection{The valuation lifting lemma}\label{subsec:lifting}

In this section, we prove that raising 
an integer $a \equiv 1 \pmod{p}$ to the $p$-th power increases its 
``distance from $1$'' in the $p$-adic sense by exactly one, provided 
the starting distance is at least $1$.

\begin{lemma}\label{lem:lifting}
Let $p$ be an odd prime, and let $a$ be an integer with $a \equiv 1 
\pmod{p}$ and $\nu_p(a - 1) = r \ge 1$. Then $\nu_p(a^p - 1) = r + 1$.
\end{lemma}

\begin{proof}
Write $a = 1 + c p^r$ with $p \nmid c$. Then
\[
  a^p - 1 = (1 + c p^r)^p - 1 = \sum_{k=1}^{p} \binom{p}{k} c^k p^{rk}.
\]
We examine the contributions term by term.

The $k = 1$ term is $\binom{p}{1} c \cdot p^r = c \cdot p^{r+1}$. 
Since $p \nmid c$, this term has $\nu_p$ equal to exactly $r + 1$.

For $2 \le k \le p - 1$, the binomial coefficient $\binom{p}{k}$ is 
divisible by $p$ (a standard fact), contributing a factor of $p$, and 
$p^{rk}$ contributes $p^{rk}$, so the term has $\nu_p \ge 1 + rk$. 
Since $r \ge 1$ and $k \ge 2$, we have $1 + rk \ge 1 + 2r \ge r + 2 > 
r + 1$.

For $k = p$, the term is $c^p p^{rp}$, which has $\nu_p \ge rp$. Since 
$p \ge 3$ (we assumed $p$ odd) and $r \ge 1$, $rp \ge 3 > r + 1$ when 
$r = 1$, and more generally $rp \ge r + 2 > r + 1$ for all $r \ge 1$.

So the $k = 1$ term has the uniquely lowest $p$-adic valuation among 
all terms; all other terms have strictly higher $\nu_p$. Therefore 
$\nu_p(a^p - 1) = r + 1$.
\end{proof}

The hypothesis $r \ge 1$ is essential: the argument collapses at 
$r = 0$, and in fact the conclusion fails there. We will apply the 
lemma only when $r \ge 1$, which is automatic from the 
definition $s = \nu_p(10^\lambda - 1) - 1 \ge 0$, so that 
$\nu_p(10^\lambda - 1) = 1 + s \ge 1$.

\subsection{Finding the exact $e$}\label{subsec:pinning}

We now apply Lemma~\ref{lem:lifting} iteratively to compute 
$\nu_p(10^{\lambda p^j} - 1)$ for all $j \ge 0$. Combined with 
Lemma~\ref{lem:form-lambda-pe}, this determines $\lambda(p^n)$ exactly.

\begin{proof}[Proof of Theorem~\ref{thm:prime-power-period}]
We claim that for all $j \ge 0$,
\begin{equation}\label{eq:valuation-growth}
  \nu_p\bigl(10^{\lambda p^j} - 1\bigr) = 1 + s + j.
\end{equation}
The case $j = 0$ is the definition of $s$. For the inductive step, 
assume \eqref{eq:valuation-growth} holds for some $j \ge 0$. Setting 
$a = 10^{\lambda p^j}$, we have $a \equiv 1 \pmod{p}$ and 
$\nu_p(a - 1) = 1 + s + j \ge 1$, so Lemma~\ref{lem:lifting} applies 
and gives $\nu_p(a^p - 1) = 2 + s + j$. Since 
$a^p = 10^{\lambda p^{j+1}}$, this is the case $j+1$ of 
\eqref{eq:valuation-growth}.

By Lemma~\ref{lem:form-lambda-pe}, $\lambda(p^n) = \lambda \cdot p^e$ 
for some $0 \le e \le n - 1$. The definition of $\lambda(p^n)$ 
requires $10^{\lambda p^e} \equiv 1 \pmod{p^n}$, which by 
\eqref{eq:valuation-growth} is equivalent to $1 + s + e \ge n$ and
$e \ge n - 1 - s$. Minimality of $\lambda(p^n)$ forces $e$ to be the 
smallest non-negative integer satisfying this, namely
\[
  e = \max(n - 1 - s,\; 0).
\]
Therefore $\lambda(p^n) = \lambda \cdot p^{\max(n-1-s,\, 0)}$.
\end{proof}


\subsection{Examples and exceptional primes}\label{subsec:examples}

We illustrate the theorem with three examples: a generic case where 
$s = 0$, and the two exceptional primes below $1000$ where $s = 1$.

\begin{example}[$p = 7$, generic case]\label{ex:p-7}
Here $\lambda = \ord_7(10) = 6$, and $\nu_7(10^6 - 1) = \nu_7(999999) 
= 1$, so $s = 0$. Theorem~\ref{thm:prime-power-period} predicts 
$\lambda(7^n) = 6 \cdot 7^{n-1}$. We verify the first few:
\[
  \lambda(7) = 6, \qquad \lambda(49) = 42, \qquad \lambda(343) = 294,
\]
matching $6 \cdot 7^0$, $6 \cdot 7^1$, $6 \cdot 7^2$ respectively.
\end{example}

\begin{example}[$p = 3$, exceptional case]\label{ex:p-3}
Here $\lambda = \ord_3(10) = 1$ (since $10 \equiv 1 \pmod{3}$), and 
$\nu_3(10^1 - 1) = \nu_3(9) = 2$, so $s = 1$. 
Theorem~\ref{thm:prime-power-period} predicts $\lambda(3^n) = 
1 \cdot 3^{\max(n - 2, 0)} = 3^{\max(n-2, 0)}$. The period stays at 
$1$ for $n = 1$ and $n = 2$, then starts growing:
\[
  \lambda(3) = 1, \quad \lambda(9) = 1, \quad \lambda(27) = 3, \quad 
  \lambda(81) = 9, \quad \lambda(243) = 27.
\]
This matches the formula exactly.
\end{example}

A search over all primes $p < 1000$ with $p \nmid 10$ found only 
$p = 3$ and $p = 487$ with $\nu_p(10^\lambda - 1) \ne 1$; all others 
have $s = 0$ and so obey the ``generic'' formula $\lambda(p^n) = 
\lambda \cdot p^{n-1}$. Primes with $s \ge 1$ are known as 
\emph{Wieferich-like primes with respect to base $10$}, and are 
conjectured but not known to be infinite. The formula in 
Theorem~\ref{thm:prime-power-period} handles both the generic and 
exceptional cases uniformly, with the parameter $s$ absorbing the 
anomaly.


%%-------------------------------------------------
\section{Proof of the Main Theorem}\label{sec:main-proof}
%%-------------------------------------------------

We now assemble the results of Sections~\ref{sec:order}, 
\ref{sec:reduction}, and~\ref{sec:prime-power} into a proof of Theorem~\ref{lem:main-formula}, restated below.

\textbf{Theorem~\ref{lem:main-formula}}. \textit{
Let $a/b$ be a fraction in lowest terms with
\[
  b = 2^\alpha \cdot 5^\beta \cdot p_1^{e_1} \cdots p_r^{e_r},
\]
where the $p_i$ are distinct primes coprime to $10$. Then the decimal
expansion of $a/b$ has exactly $L = \max(\alpha, \beta)$ non-repeating
digits, followed by a purely repeating block of length
\[
  \lcm\!\bigl(\lambda(p_1^{e_1}),\, \ldots,\, \lambda(p_r^{e_r})\bigr),
\]
where each
\[
  \lambda(p_i^{e_i}) \;=\; \lambda(p_i)\cdot p_i^{\max(e_i - 1 - s_i,\; 0)},
  \qquad s_i = \nu_{p_i}(10^{\lambda(p_i)} - 1) - 1.
\]}

\begin{proof}
Write $m = p_1^{e_1} \cdots p_r^{e_r}$, so that $b = 2^\alpha \cdot 
5^\beta \cdot m$ with $\gcd(m, 10) = 1$.

\smallskip

\emph{Pre-period.} By Theorem~\ref{thm:preperiod}, the decimal 
expansion of $a/b$ has exactly $L = \max(\alpha, \beta)$ non-repeating 
digits, after which it is purely periodic with period $\ord_m(10)$.

\smallskip

\emph{Period.} By Lemma~\ref{lem:ord-lcm},
\[
  \ord_m(10) = \lcm\!\bigl(\ord_{p_1^{e_1}}(10),\, \ldots,\, 
  \ord_{p_r^{e_r}}(10)\bigr),
\]
and by Theorem~\ref{thm:prime-power-period}, each factor is
\[
  \ord_{p_i^{e_i}}(10) = \lambda(p_i^{e_i}) = \lambda(p_i) \cdot 
  p_i^{\max(e_i - 1 - s_i,\, 0)}
\]
with $s_i = \nu_{p_i}(10^{\lambda(p_i)} - 1) - 1$.

\smallskip

\emph{Independence from $a$.} By Theorem~\ref{thm:period-equals-order} 
and its corollary, $\ord_m(10)$ depends only on $m$ (and hence only on 
$b$), not on the numerator $a$. The pre-period $L = \max(\alpha, \beta)$ 
also depends only on $b$.

\smallskip

Combining these three observations yields the Main Theorem.
\end{proof}

We close with a worked example that exercises every component of the 
theorem.

\begin{example}[$13/1400$]\label{ex:1-over-1400}
Factoring the denominator, $1400 = 2^3 \cdot 5^2 \cdot 7$, so 
$\alpha = 3$, $\beta = 2$, $m = 7$, and $L = \max(3, 2) = 3$. The 
pre-period is $3$ digits. For the period, the only prime power in $m$ 
is $7$ itself, and by Example~\ref{ex:p-7} we have $\lambda(7) = 6$. 
So the period is $\lcm(6) = 6$. Therefore, Theorem~\ref{lem:main-formula} therefore predicts 
three non-repeating digits followed by a block of length $6$. Indeed,
\[
  \frac{13}{1400} = 0.009\overline{285714},
\]
confirming the prediction.
\end{example}

%%-------------------------------------------------
\section{Group Theory Perspective on Periods} \label{sec:group_theory}
%%-------------------------------------------------
Following the previous sections, we have derived a method to determine the period of an arbitrary fraction by decomposing it into a product of prime-power fractions. This allows us to find the period of any fraction $\frac{a}{b}$ given that we know the period of $\frac{1}{p}$ for all prime natural numbers $p$ that divide $b$. In this section, we will utilize the results of Lemma \ref{lem:period-as-divisibility} and analyze the group theory implications of this lemma to create conditions for the values of the possible periods of $\frac{1}{p}$. Following this, we will also look how, as a consequence of this lemma, we can determine classes of natural numbers that certain to have specific periods.\\

To begin, the condition proven in Lemma \ref{lem:period-as-divisibility}, when applied to a fraction $\frac{a}{b}=\frac{1}{p}$ reads that the period of $\frac{1}{p}$ is the smallest positive integer $k$ such that $p$ | $10^k-1$ for all $p\notin \{2,5\}$. This is equivalent to stating that $10^k=n*p+1$ for some arbitrary integer $n$. Redefining this as a statement in terms of modulo $p$, this condition is equivalent to stating that $k$ is the smallest positive integer $k$ such that $10^k=1\mod{p}$.\\

By redefining Lemma \ref{lem:period-as-divisibility} in this way, we see an alternative representation of this problem of determining the period of $\frac{1}{p}$ in a group theory context. Namely, this problem is the same as determining the order of the subgroup generated by $10$ in the multiplicative group of integers modulo $p$, a group that we will refer to as $(\mathbb{Z}/p\mathbb{Z})^X$. Using this definition that $\lambda(\frac{1}{p})=\ord_p(10)$ allows us to apply various theorems relating to cyclic groups, like $(\mathbb{Z}/p\mathbb{Z})^X$, to limit the possible values of the size of this subgroup, and in turn the period of $\frac{1}{p}$. We will now prove some of the properties of $\lambda(\frac{1}{p})$ using the group structure established and Lagrange's Theorem.

\begin{lemma} \label{lem:lagrange_divisibility}
    For any prime $p$, $\lambda(\frac{1}{p}) | (p-1)$
\end{lemma}
\begin{proof}
    By Lagrange's Theorem, since $10$ generates a subgroup of order $\ord_p(10)$ within $(\mathbb{Z}/p\mathbb{Z})^X$, the order of that subgroup must divide the order of $(\mathbb{Z}/p\mathbb{Z})^X$. The order of the multiplicative group of integers modulo $p$ is given by $\phi(p)$, the Euler totient function of $p$. However, since $p$ is prime $\phi(p)=p-1$, so $(\mathbb{Z}/p\mathbb{Z})^X$ has order $p-1$. As a result, it must be the case that $\lambda(\frac{1}{p})=\ord_p(10)|(p-1)$. Giving the desired conclusion that $$\lambda(\frac{1}{p}) | (p-1)$$
\end{proof}

This result allows us to drastically limit the possible values of $\lambda(\frac{1}{p})$ from essentially any positive integer being a candidate, to only the select few integers that divide $p-1$. This result also corroborates the proof that periods cannot be infinite, as $\ord_p(10)$ must be less than or equal to $p-1$ for Lagrange's Theorem to hold and in turn $\frac{1}{p}$ cannot have a period as large as or larger than its denominator. This places an upper bound on the period of $\frac{1}{p}$. Now, within this framework, we can also form a lower bound, proving that primes of a sufficient size can only have a certain range of period lengths.

\begin{theorem} \label{thm:period_bounds}
    $\log(p) < \lambda(\frac{1}{p})\leq p-1$
\end{theorem}
\begin{proof}
    \underline{Upper Bound}: Follows directly from Lemma \ref{lem:lagrange_divisibility}. Since $\lambda(\frac{1}{p}) | (p-1)$, it follows that $\lambda(\frac{1}{p})\leq p-1$ \\
    \underline{Lower Bound}: For the sake of contradiction, assume that $\log(p) \geq \lambda(\frac{1}{p})$, then that implies that $10^{\log(p)} \geq 10^{\lambda(\frac{1}{p})}=10^k$. Rearranging this inequality, we get that $10^k\leq p$. However, if $10^k\leq p$, then since $10^k$ is always divisible by 10, it cannot be prime, and so $ p\neq 10^k$. This means that the only possible case is that $10^k< p$ in which case $10^k=10^k\mod{p}$. However, since $k$ cannot equal $0$, $10^k\neq1\mod{p}$ which contradicts the definition of $k$. As such, it's impossible for $\log(p) \geq \lambda(\frac{1}{p})$, it must be that $\log(p) < \lambda(\frac{1}{p})$.\\
    
    Combining the two inequalities, we get the desired inequality that $$\log(p) < \lambda(\frac{1}{p})\leq p-1$$
\end{proof}

While this is a loose bound for large primes, it allows for us to eliminate certain potential period lengths on the lower end of the spectrum. For example, for any $p>100$, the lower bound establishes that even though $p-1$ must be even, and is therefore divisible by both $1$ and $2$, neither may be the period of $\frac{1}{p}$ since $log(p)>2$. As such, the lower bound allows for the elimination of many of the common small factors of $p-1$ as candidates for $\lambda(\frac{1}{p})$ once $p$ is sufficiently large. Another consequence of this is that the period of $\frac{1}{p}$ is generally increasing as $p$ gets larger as this lower bound increases the minimum value of the period.\\

Having established an upper and lower bound for $\lambda(\frac{1}{p})$, along with the condition of needing to divide $p-1$, we have greatly reduced the possible values of the period for an arbitrary fraction $\frac{1}{p}$. Now that we have created a method to determine the period of an arbitrary fraction, we ask the opposite. Is it possible to find a fraction of an arbitrary period? To answer this question, we will look back on Lemma \ref{lem:period-as-divisibility} and see how by interpreting this condition differently, we can map out all fractions of the form $\frac{1}{b}$ with a desired period and determine a set of elements for which we can immediately determine their period just by inspection.\\

\subsection{Determining all unit fractions of period $k$}
Earlier in this section, we mentioned that another representation of Lemma \ref{lem:period-as-divisibility} is that $10^k=n*b+1$ for some arbitrary integer $n$. As the original lemma only mentions that $b$ must be coprime to 10, this condition also holds for not necessarily prime integers $b$, so long as $b$ does not have factors of $2$ or $5$. From this representation, we get that $\frac{1}{b}$ has period $k$ if and only if $k$ is the smallest positive integer, such that $(10^k-1)$ is divisible by $b$. We can use this definition to map out ever possible fraction of period $k$ recursively by taking $k=\gamma$, finding the prime factorization of $10^\gamma-1$ and then finding all combinations of prime factors that produce integers, $\tau_i$, not within the range of fractions of any shorter period. This will mean that $k=\gamma$ is the smallest case where the condition holds, and as a result, $\frac{1}{\tau_i}$ must have period $\gamma$ for all $i$. We will start with the $k=1$, $k=2$, and $k=3$ case, and then mention a pattern that can be noticed for any general value of $k$.

\begin{example}[Unit fractions of period $k=1$] Set $k=1$, we get $10^1-1=n*b$ which becomes $9=n*b$. Now decomposing $9=3*3*1$, we find that the only values that can satisfy $9=n*b$ are the arbitrary products of this decomposition. As such, ignoring $1$, we get the values $\{3,9\}$ are the only possible values of $b$ for $k=1$. As such $\lambda(\frac{1}{3})=\lambda(\frac{1}{9})=1$. However, as shown in Theorem \ref{thm:preperiod}, an additional factor of $\frac{1}{2^\alpha \cdot5^\beta}$ only increases the number of non-repeating decimals without changing the period, we can represent all fractions with period $k=1$ as $\frac{1}{3\cdot2^\alpha \cdot5^\beta}$ or $\frac{1}{9\cdot2^\alpha \cdot5^\beta}$ for any natural number $\alpha,\beta$ including $0$.
\end{example}

\begin{example}[Unit fractions of period $k=2$] Now, set $k=2$, we get $10^2-1=n*b$ which becomes $99=n*b$, where $99=3*3*11$. All of the products of these prime factors are $\{3,9,11,33,99\}$. However, since we are searching for the minimum value of $k$ such that these factors divide $10^k-1$, we must remove the factors that appeared for $k=1$. As such, the list decreases to $\{11,33,99\}$. Similarly, we now know that all unit fractions with period $\lambda=2$ can be represented as one of $\{\frac{1}{11\cdot2^\alpha \cdot5^\beta},\frac{1}{33\cdot2^\alpha \cdot5^\beta},\frac{1}{99\cdot2^\alpha \cdot5^\beta}\}$
\end{example}

\begin{example}[Unit fractions of period $k=3$] For $k=3, 10^3-1=n*b$, so $n*b=999=3*3*111=3*3*3*37$. With these factors, we get that the new products of these factors are $\{27,111,333,999\}$. So all of the fractions of period $k=3$ can be written as $\{\frac{1}{27\cdot2^\alpha \cdot5^\beta},\frac{1}{111\cdot2^\alpha \cdot5^\beta},\frac{1}{333\cdot2^\alpha \cdot5^\beta},\frac{1}{999\cdot2^\alpha \cdot5^\beta}\}$
\end{example}

As we see from these examples, we are capable of determining every unit fraction of a given period by following this recursive method of taking the prime factorization of $10^k-1$ and finding the combination of products that have not appeared in any smaller values of $k$. However, this process still makes it difficult to find a fraction of a large value of $k$, like for example $k=493$. However, we can notice a pattern about this recursive process that allows us to identify infinitely many, though not necessarily all, unit fractions of any arbitrary period $k$ without the need for this recursive process. Specifically, we must look at what the value $10^k-1$ for an arbitrary $k$.\\

By inspection, we see that taking $10^k-1$ for any arbitrary $k$ with give use a repetition of $9$ repeated $k$ times. Since all of the $9$s are divisible by both $3$ and $9$, we get that $10^k-1=3*3*\frac{10^k-1}{9}$, where since we've divided the repeated nines by $3*3=9$, we get that $\frac{10^k-1}{9}$ is equal to $1$ repeated $k$ times. These terms of repeated ones are known as a repunit, and we will denote the $k$th repunit as $R_k=\frac{10^k-1}{9}$. While we can't tell if $R_k$ can be further decomposed, we can now prove an important statement about these repunits. That is, that $\lambda(\frac{1}{R_k})=k$ for any $k>1$  
\begin{theorem} \label{thm:Repunit_period}
    $\lambda(\frac{1}{R_k})=k$ for all $k>1$ within the positive integers
\end{theorem}

\begin{proof}
    To show that $\lambda(\frac{1}{R_k})=k$, we must show that $R_k$ divides $10^k-1$, and also that $R_k$ does not divide $10^{k'}-1$ for any $k'<k$. \\
    The first point comes by the definition of the repunits. Since $R_k=\frac{10^k-1}{9}$, $\frac{10^k-1}{R_k}=9$ for all $k$. This already tells us that $\lambda(\frac{1}{R_k})\leq k$ for all $k$, so to prove that they are equal, we need to show that there exists no $k'<k$ such that $R^k | 10^{k'-1}$.\\
    To do this, we can look at the structure of the repunits. By definition $R_{k'}=\frac{10^{k'}-1}{9}<10^{k'} $, where $10^k$ is equal to $1$ followed by $k'$ zeroes. While $R_k$ is equal to one repeated $k$ times, which is equivalent to one repeated $k$ times or 1 followed by $k-1$ ones. However if $k>k'$, then $k-1\geq k'$. If $k-1>k'$, then $R_k>10^{k'}$ since $R_k$ has more digits, and if $k-1=k'$, $R_k$ is still larger than $10^k$ since while the two may have the same number of digits, the repeated ones are larger than the repeated zeroes. So, in either case $R_k>10^{k'}$ and as a result, cannot divide $10^{k'}-1$ for any $k'<k$. \\

    Since $R_k$ divides $10^k-1$ and does not divide $10^{k'}-1$ for any $k'<k$, it must be the case the $\frac{1}{R_k}$ has period $k$
\end{proof}
From this realization, we can create infinitely unit fractions with period $k$ by taking the $k$th repunit as using the fact the a $2^\alpha*5^\beta$ term does not affect the period, tho get that there are infinitely many unit fractions of the form $\{\frac{1}{R_k*2^\alpha*5^\beta},\frac{1}{3*R_k*2^\alpha*5^\beta},\frac{1}{9*R_k*2^\alpha*5^\beta}\}$ for any positive integer $\alpha,\beta$ including $0$. The times $\frac{1}{3}$ and times $\frac{1}{9}$ follows from the fact that the same logic can be used to show that $\lambda(\frac{1}{R_k})=\lambda(\frac{1}{3R_k})=\lambda(\frac{1}{9R_k})=k$\\

While these are not necessarily all of the unit fractions of period $k$, they can all be found using the recursive method outlined in the previous examples. However, knowing that $\frac{1}{R_k}$ has period $k$ for all $k>1$, and that there are repunits of any positive integer values of $k$ allows us to find a fraction of any arbitrary period we desire. As such we have outlined a recursive method to determine all of the unit fractions of period $k$, while also determining that $\frac{1}{R_k}$ is guaranteed to have period $k$, for all $k>1$, allowing us to easily find a selection of fractions of arbitrary period.
%%-------------------------------------------------
\section{A Probabilistic Perspective on Periods}\label{sec:probabilistic}
%%-------------------------------------------------

Sections~\ref{sec:order}--\ref{sec:main-proof} give a complete formula for
the period of any single fraction $a/b$. A natural next question is how
those periods are distributed as the prime varies: of all primes
$p \le X$, what fraction yield a period equal to $p - 1$? What fraction
yield a period equal to $(p-1)/2$? In this section we recast the period
as the index of $10$ in the cyclic group $(\Z/p\Z)^*$, derive the exact
distribution of that index when one samples a uniformly random element of
the group, and use it to predict the empirical distribution of periods
across primes. The case of maximal period turns out to be classical: it is
exactly the setting of Artin's conjecture.

\subsection{The index of an element and its connection to $\lambda(p)$}
\label{subsec:index}

Throughout this section, $p$ denotes a prime with $p \neq 2$ and $p \neq 5$,
and $G = (\Z/p\Z)^*$ denotes the multiplicative group of nonzero residues
modulo $p$, which is cyclic of order $p - 1$. As before, $\lambda(p) =
\ord_p(10)$, and $\varphi$ is Euler's totient function.

\begin{definition}\label{def:index}
For $g \in G$, the \emph{index} of $g$ is the positive integer
\[
    \ind(g) \;:=\; \frac{|G|}{\ord(g)} \;=\; \frac{p - 1}{\ord(g)}.
\]
\end{definition}

By Lagrange's theorem applied to the cyclic subgroup $\langle g \rangle$,
$\ord(g)$ divides $|G| = p - 1$, so $\ind(g)$ is always a positive divisor
of $p - 1$. The index is the standard group-theoretic index of
$\langle g \rangle$ inside $G$.

Specializing to $g = 10$ and combining with
Theorem~\ref{thm:period-equals-order},
\[
    \ind(10) \;=\; \frac{p - 1}{\lambda(p)},
\]
so the index of $10$ records exactly how the period of $1/p$ compares to
the maximum possible value $p - 1$. Index $1$ corresponds to a
\emph{full-reptend prime} (period $p - 1$); index $d$ in general
corresponds to period $(p - 1)/d$.

\subsection{Exact distribution of the index in $(\Z/p\Z)^*$}
\label{subsec:distribution}

For each prime $p$, take the probability space $\Omega_p = G$ with the
uniform measure $\mathbb{P}(\{g\}) = 1/(p-1)$, and let $D_p$ denote the
index of a uniformly random element of $G$. The following theorem gives
the distribution of $D_p$ in closed form.

\begin{theorem}\label{thm:index-distribution}
Let $p$ be a prime with $p \neq 2$ and $p \neq 5$, and let $d$ be a
positive divisor of $p - 1$. Then
\[
    \mathbb{P}(D_p = d) \;=\; \frac{\varphi\!\left((p-1)/d\right)}{p - 1}.
\]
If $d \nmid (p-1)$, then $\mathbb{P}(D_p = d) = 0$.
\end{theorem}

\begin{proof}
Since $G$ is cyclic of order $p - 1$, fix a generator $g_0$. Every element
of $G$ can be written uniquely as $g_0^k$ for some $k \in \{0, 1, \ldots,
p - 2\}$. By Lagrange's theorem applied to cyclic groups,
$\ord(g_0^k) = (p - 1)/\gcd(k, p - 1)$, so
\[
    D_p(g_0^k) = d
    \;\Longleftrightarrow\;
    \ord(g_0^k) = \frac{p - 1}{d}
    \;\Longleftrightarrow\;
    \gcd(k, p - 1) = d.
\]
The condition $\gcd(k, p - 1) = d$ forces $d \mid k$, so write $k = d\ell$
with $\ell \in \{0, 1, \ldots, (p-1)/d - 1\}$. Factoring $d$ out of both
arguments of the gcd,
\[
    \gcd(k, p - 1) \;=\; \gcd\!\left(d \ell,\, d \cdot \tfrac{p-1}{d}\right)
    \;=\; d \cdot \gcd\!\left(\ell,\, \tfrac{p-1}{d}\right),
\]
so $\gcd(k, p - 1) = d$ if and only if $\gcd(\ell, (p-1)/d) = 1$. By the
definition of Euler's totient, the number of such $\ell$ is
$\varphi((p-1)/d)$. Dividing by $|G| = p - 1$ yields the stated
probability. If $d \nmid (p-1)$, no element of $G$ has order $(p-1)/d$,
so the probability is $0$.
\end{proof}

\begin{example}\label{ex:p-7-distribution}
For $p = 7$ we have $p - 1 = 6$, with divisors $1, 2, 3, 6$.
Theorem~\ref{thm:index-distribution} gives
\[
    \mathbb{P}(D_7 = 1) = \tfrac{\varphi(6)}{6} = \tfrac{2}{6}, \quad
    \mathbb{P}(D_7 = 2) = \tfrac{\varphi(3)}{6} = \tfrac{2}{6}, \quad
    \mathbb{P}(D_7 = 3) = \tfrac{\varphi(2)}{6} = \tfrac{1}{6}, \quad
    \mathbb{P}(D_7 = 6) = \tfrac{\varphi(1)}{6} = \tfrac{1}{6},
\]
summing to $1$ as expected.
\end{example}

\subsection{Averaging over primes: a conjecture}\label{subsec:averaging}

Theorem~\ref{thm:index-distribution} is a statement about a single prime
$p$. To predict the distribution of periods $\lambda(p)$ as $p$ varies,
we average the per-prime probability across primes.

Let
\[
    \pi(X) \;:=\; |\{p \le X : p \text{ prime},\; p \neq 2,\; p \neq 5\}|,
\]
and let $p_d(X)$ denote the observed proportion of such primes whose
period index $\ind(10) = (p-1)/\lambda(p)$ equals $d$. Define the
\emph{predicted proportion} by averaging the per-prime probability from
Theorem~\ref{thm:index-distribution}:
\[
    \widehat{p}_d(X)
    \;:=\; \frac{1}{\pi(X)}
    \sum_{\substack{p \le X \\ p \text{ prime},\; p \neq 2,\; p \neq 5 \\
    d \mid p - 1}} \frac{\varphi((p-1)/d)}{p - 1}.
\]
Note that $\widehat{p}_d(X)$ is purely a cyclic-group quantity: it
records, averaged over primes $p \le X$, the probability that a uniformly
random element of $(\Z/p\Z)^*$ has index $d$. It does not refer to $10$
in any way.

\begin{conjecture}\label{conj:agreement}
For every fixed positive integer $d$, both limits
\[
    p_d^\infty \;:=\; \lim_{X \to \infty} p_d(X), \qquad
    \widehat{p}_d^{\,\infty} \;:=\; \lim_{X \to \infty} \widehat{p}_d(X)
\]
exist, and they are equal.
\end{conjecture}

We computed both quantities at $X = 10{,}000$ for the smallest values of
$d$, and they agree to within $1$--$4\%$:

\begin{center}
\begin{tabular}{rrrr}
\hline\hline
Index $d$ & Observed $p_d(10{,}000)$ & Predicted $\widehat{p}_d(10{,}000)$
& Ratio \\
\hline
1 & 38.1\% & 37.4\% & 1.02 \\
2 & 29.3\% & 28.2\% & 1.04 \\
3 & 6.9\%  & 6.6\%  & 1.03 \\
4 & 6.9\%  & 7.0\%  & 0.98 \\
6 & 5.0\%  & 5.0\%  & 0.99 \\
8 & 1.7\%  & 1.7\%  & 1.00 \\
\hline\hline
\end{tabular}
\end{center}

The closeness of the two columns is strong empirical evidence for
Conjecture~\ref{conj:agreement} in the range we can compute: the period
of $1/p$ behaves, statistically, exactly as if $10$ were a uniformly
random element of $(\Z/p\Z)^*$.

\subsection{Connection to Artin's conjecture}\label{subsec:artin}

The case $d = 1$ of Conjecture~\ref{conj:agreement} has a classical
interpretation. The condition $\ind(10) = 1$ means $\langle 10 \rangle =
(\Z/p\Z)^*$, i.e. $10$ is a primitive root modulo $p$, equivalently
$\lambda(p) = p - 1$. Such primes are the \emph{full-reptend primes}
for base $10$. Artin's conjecture (1927) predicts that for any fixed
non-square integer base, a positive proportion of primes have that base
as a primitive root; for base $10$, the predicted density is Artin's
constant, approximately $0.3739$. Our computation in
Section~\ref{subsec:averaging} gives $\widehat{p}_1(10{,}000) \approx
0.38$, in striking agreement with this prediction, although Artin's
conjecture itself remains open.
\section{Acknowledgements}
Sky Pulling is responsible for Sections 1-4\\
Jovani Pitterson is responsible for Section 5\\
Samir Kadariya is responsible for the Introduction and Section 6
\end{document}